\chapter{}
\section{Gruppentheorie}

\dfn{Gruppe}{Eine \textbf{Gruppe} \((G,\circ)\) besteht aus einer nichtleeren Menge \(G\) und einer Verknüpfung
\[
\circ\colon G\times G\to G,
\]
die folgende Axiome erfüllt:
\begin{enumerate}[label=\alph*)]
\item \textbf{Assoziativität:} Für alle \(a,b,c\in G\) gilt
\[
(a\circ b)\circ c=a\circ(b\circ c).
\]

\item \textbf{Neutrales Element:} Es gibt ein Element \(e\in G\), sodass für alle \(g\in G\)
\[
e\circ g=g\circ e=g.
\]

\item \textbf{Inverse Elemente:} Für jedes \(g\in G\) gibt es ein \(h\in G\), sodass
\[
g\circ h=h\circ g=e.
\]
Das Element \(h\) heißt Inverses zu \(g\) und wird mit \(g^{-1}\) bezeichnet.
\end{enumerate}
Gilt zusätzlich
\[
g\circ h=h\circ g\qquad\text{für alle }g,h\in G,
\]
so heißt die Gruppe \textbf{abelsch} oder \textbf{kommutativ}.}

\ex{Grundlegende Beispiele}{
\begin{enumerate}[label=\alph*)]
\item \((\bbZ,+)\) ist eine abelsche Gruppe. Das neutrale Element ist \(0\), denn \(a+0=0+a=a\). Das Inverse von \(a\) ist \(-a\), denn \(a+(-a)=(-a)+a=0\).

\item \((\bbR\setminus\{0\},\cdot)\) ist eine abelsche Gruppe. Das neutrale Element ist \(1\), und das Inverse eines Elements \(a\neq0\) ist \(a^{-1}=1/a\). Das Produkt zweier von null verschiedener reeller Zahlen ist wieder von null verschieden, sodass die Verknüpfung tatsächlich auf dieser Menge definiert ist.

\item Sei \(M\) eine nichtleere Menge und
\[
\operatorname{Bij}(M):=\{f\colon M\to M\mid f\text{ bijektiv}\}.
\]
Dann ist \((\operatorname{Bij}(M),\circ)\) eine Gruppe. Die Verknüpfung ist die Komposition von Abbildungen, das neutrale Element ist die Identität \(\operatorname{id}_M\), und das Inverse von \(f\) ist die Umkehrabbildung \(f^{-1}\). Die Assoziativität folgt aus der Assoziativität der Komposition.

\item \((\bbN,+)\) ist keine Gruppe: Es gibt kein neutrales Element in \(\bbN=\{1,2,\ldots\}\). Auch \((\bbN_0,+)\) ist keine Gruppe, denn zwar ist \(0\) neutral, aber beispielsweise besitzt \(1\) kein additives Inverses in \(\bbN_0\).

\item Sind \((G_1,\circ)\) und \((G_2,\ast)\) Gruppen, so wird \(G_1\times G_2\) mit der komponentenweisen Verknüpfung
\[
(g_1,g_2)\cdot(h_1,h_2)
=(g_1\circ h_1,g_2\ast h_2)
\]
ebenfalls zu einer Gruppe. Das neutrale Element ist \((e_1,e_2)\), und das Inverse von \((g_1,g_2)\) ist \((g_1^{-1},g_2^{-1})\).
\end{enumerate}}

\dfn{Addition modulo \(m\)}{Sei \(m\in\bbN\). Wir betrachten die Äquivalenzklassen der Kongruenzrelation modulo \(m\) auf \(\bbZ\). Die Menge dieser Klassen wird mit
\[
\bbZ/m\bbZ
\]
bezeichnet. Auf ihr definieren wir
\[
[x]+[y]:=[x+y].
\]
Diese Verknüpfung ist wohldefiniert: Das Ergebnis hängt nicht von der Wahl der Vertreter \(x\) und \(y\) ab. Mit dieser Verknüpfung ist \(\bbZ/m\bbZ\) eine abelsche Gruppe.}

\nt{Die Wohldefiniertheit ist wichtig, weil eine Äquivalenzklasse viele verschiedene Vertreter haben kann. Wenn wir etwa \([x]=[x']\) und \([y]=[y']\) haben, müssen wir sicherstellen, dass auch
\[
[x+y]=[x'+y']
\]
gilt. Dies folgt daraus, dass \(m\mid(x-x')\) und \(m\mid(y-y')\) zusammen \(m\mid((x+y)-(x'+y'))\) ergeben.}

\mlemma{Eindeutigkeit des neutralen und inversen Elements}{Sei \((G,\circ)\) eine Gruppe.
\begin{enumerate}[label=\alph*)]
\item Das neutrale Element \(e\) ist eindeutig bestimmt.
\item Zu jedem \(a\in G\) ist das inverse Element \(a^{-1}\) eindeutig bestimmt.
\end{enumerate}}

\begin{myproof}
\textbf{a)} Seien \(e_1\) und \(e_2\) zwei neutrale Elemente. Dann gilt
\[
e_1=e_1\circ e_2=e_2,
\]
weil \(e_2\) rechtsneutral und \(e_1\) linksneutral ist. Also sind sie gleich.

\textbf{b)} Seien \(b_1\) und \(b_2\) beide inverse Elemente zu \(a\). Dann gilt
\[
\begin{aligned}
b_1
&=b_1\circ e\\
&=b_1\circ(a\circ b_2)\\
&=(b_1\circ a)\circ b_2\\
&=e\circ b_2\\
&=b_2.
\end{aligned}
\]
Also ist das inverse Element eindeutig.
\end{myproof}

\thm{Rechenregeln in Gruppen}{Sei \((G,\circ)\) eine Gruppe und seien \(a,b,x,x'\in G\).
\begin{enumerate}[label=\alph*)]
\item Das Inverse eines Produkts ist
\[
(a\circ b)^{-1}=b^{-1}\circ a^{-1}.
\]

\item Es gelten die Kürzungsregeln:
\[
a\circ x=a\circ x'\implies x=x',
\]
sowie
\[
x\circ a=x'\circ a\implies x=x'.
\]

\item Für jedes \(a\in G\) sind die Abbildungen
\[
r_a\colon G\to G,\quad g\mapsto g\circ a,
\]
und
\[
l_a\colon G\to G,\quad g\mapsto a\circ g
\]
bijektiv. Sie heißen Rechts- beziehungsweise Linkstranslation.
\end{enumerate}}

\begin{myproof}
\textbf{a)} Wir zeigen, dass \(b^{-1}\circ a^{-1}\) das Inverse von \(a\circ b\) ist:
\[
\begin{aligned}
(a\circ b)\circ(b^{-1}\circ a^{-1})
&=a\circ(b\circ b^{-1})\circ a^{-1}\\
&=a\circ e\circ a^{-1}=e.
\end{aligned}
\]
In umgekehrter Reihenfolge ergibt sich ebenso
\[
(b^{-1}\circ a^{-1})\circ(a\circ b)=e.
\]
Wegen der Eindeutigkeit des Inversen folgt die Behauptung.

\textbf{b)} Aus \(a\circ x=a\circ x'\) folgt durch Verknüpfung von links mit \(a^{-1}\):
\[
x=e\circ x=a^{-1}\circ a\circ x
=a^{-1}\circ a\circ x'=e\circ x'=x'.
\]
Die zweite Kürzungsregel wird analog durch Verknüpfung von rechts mit \(a^{-1}\) bewiesen.

\textbf{c)} Die Abbildung \(r_a\) besitzt die Umkehrabbildung
\[
r_{a^{-1}}\colon g\mapsto g\circ a^{-1},
\]
denn \(r_{a^{-1}}(r_a(g))=(g\circ a)\circ a^{-1}=g\). Analog ist \(l_{a^{-1}}\) die Umkehrabbildung von \(l_a\). Damit sind beide Abbildungen bijektiv.
\end{myproof}

\cor{Translationsinvarianz der Mächtigkeit}{Sei \(G\) eine Gruppe, \(M\subseteq G\) endlich und \(g\in G\). Dann gilt
\[
|gM|=|M|,
\qquad
gM:=\{g\circ m\mid m\in M\}.
\]
Denn die Linkstranslation \(m\mapsto g\circ m\) ist bijektiv und erhält daher die Anzahl der Elemente.}

\ex{Die Gruppentafel von \(\bbZ/3\bbZ\)}{
Die Gruppe \((\bbZ/3\bbZ,+)\) besteht aus den drei Klassen
\([0]\), \([1]\) und \([2]\). Die Addition erfolgt modulo \(3\):
\[
\begin{array}{c|ccc}
+ &[0]&[1]&[2] \\ \hline
{[0]}&[0]&[1]&[2] \\
{[1]}&[1]&[2]&[0] \\
{[2]}&[2]&[0]&[1]
\end{array}
\]
Beispielsweise gilt $[1]+[2]=[0]$, da $1+2=3$
durch $3$ teilbar ist. Das neutrale Element ist $[0]$.
Die Elemente $[1]$ und $[2]$ sind zueinander invers.
}

\nt{Für \(g\in G\) und \(n\in\bbN\) schreiben wir
\[
g^n=\underbrace{g\circ\cdots\circ g}_{n\text{-mal}},
\qquad
g^{-n}=\underbrace{g^{-1}\circ\cdots\circ g^{-1}}_{n\text{-mal}}.
\]
Außerdem setzen wir \(g^0=e\). Für \(n,m\in\bbZ\) gilt dann
\[
g^n\circ g^m=g^{n+m}.
\]
Diese Schreibweise beschreibt wiederholte Verknüpfung in der Gruppe. Sie muss nicht mit gewöhnlichen Potenzen reeller Zahlen verwechselt werden. Häufig lassen wir das Verknüpfungszeichen weg und schreiben \(gh\) statt \(g\circ h\).}

\dfn{Ordnung einer Gruppe und eines Elements}{Sei \(G\) eine Gruppe.
\begin{enumerate}[label=\alph*)]
\item Ist \(G\) endlich, so heißt \(|G|\) die \textbf{Ordnung der Gruppe}.
\item Sei \(g\in G\). Falls es ein kleinstes \(n\in\bbN\) mit \(g^n=e\) gibt, heißt dieses \(n\) die \textbf{Ordnung des Elements} \(g\).
\end{enumerate}}

\dfn{Symmetrische Gruppe}{Sei \(n\in\bbN\). Die Menge aller bijektiven Abbildungen
\[
\pi\colon\{1,\ldots,n\}\to\{1,\ldots,n\}
\]
heißt \textbf{symmetrische Gruppe} und wird mit \(S_n\) bezeichnet. Die Verknüpfung ist die Komposition von Abbildungen. Ihre Elemente heißen \textbf{Permutationen}.}

\nt{Eine Permutation ordnet die Zahlen \(1,\ldots,n\) neu an. Weil sie bijektiv ist, kommt jede Zahl genau einmal als Bild vor. Für \(\pi\in S_n\) gibt es \(n\) Möglichkeiten für \(\pi(1)\), danach noch \(n-1\) Möglichkeiten für \(\pi(2)\), und so weiter. Daher gilt
\[
|S_n|=n!.
\]}

\ex{Die Gruppe \(S_3\)}{
Die Gruppe \(S_3\) besteht aus den sechs Permutationen der Menge \(\{1,2,3\}\). Beispielsweise ist
\[
\pi=
\begin{pmatrix}
1&2&3\\
2&3&1
\end{pmatrix}
\]
die Abbildung \(1\mapsto2\), \(2\mapsto3\), \(3\mapsto1\).

Die Gruppe \(S_n\) ist für \(n\geq3\) nicht abelsch. Im Allgemeinen hängt die Komposition zweier Permutationen also von der Reihenfolge ab:
\[
\pi\circ\sigma\neq\sigma\circ\pi.
\]
}

\dfn{Zyklen}{Eine Permutation \(\pi\in S_n\) heißt \textbf{Zyklus}, wenn es paarweise verschiedene Elemente \(a_1,\ldots,a_m\) gibt, sodass
\[
\pi(a_i)=a_{i+1}\quad(1\leq i<m),
\qquad
\pi(a_m)=a_1,
\]
und alle übrigen Elemente fest bleiben. Man schreibt
\[
\pi=(a_1\,a_2\,\cdots\,a_m).
\]
Ein Zyklus der Länge \(2\) heißt \textbf{Transposition}. Zwei Zyklen heißen \textbf{disjunkt}, wenn sie keine gemeinsamen Elemente bewegen. Disjunkte Zyklen kommutieren.}

\ex{Zyklenschreibweise}{
Der Zyklus \((1\,2\,3)\) bildet \(1\) auf \(2\), \(2\) auf \(3\) und \(3\) auf \(1\) ab. Alle anderen Zahlen bleiben unverändert. Die Schreibweisen
\[
(1\,2\,3)=(2\,3\,1)=(3\,1\,2)
\]
beschreiben dieselbe Permutation, weil sie dieselbe zyklische Reihenfolge angeben.
}

\thm{Zerlegung in disjunkte Zyklen}{Jede Permutation \(\pi\in S_n\) lässt sich als Produkt paarweise disjunkter Zyklen schreiben.}

\begin{myproof}
Wähle ein Element \(a_1\in\{1,\ldots,n\}\). Betrachte die Folge
\[
a_1,\ \pi(a_1),\ \pi^2(a_1),\ \ldots.
\]
Da die Menge endlich ist, müssen zwei Elemente der Folge übereinstimmen. Weil \(\pi\) bijektiv ist, folgt daraus, dass es ein kleinstes \(m\geq1\) gibt, für das \(\pi^m(a_1)=a_1\) gilt. Die Elemente
\[
a_1,\pi(a_1),\ldots,\pi^{m-1}(a_1)
\]
sind paarweise verschieden und bilden einen Zyklus.

Falls dieser Zyklus noch nicht alle Elemente umfasst, wählen wir ein Element \(a_2\), das darin nicht vorkommt, und konstruieren auf dieselbe Weise einen weiteren Zyklus. Die so entstehenden Zyklen sind disjunkt: Würde die Bahn eines Elements aus dem ersten Zyklus mit der Bahn eines Elements aus dem zweiten zusammentreffen, lägen beide Elemente in derselben Bahn und damit im selben Zyklus.

Da die Menge endlich ist, endet dieses Verfahren nach endlich vielen Schritten. Das Produkt der entstandenen disjunkten Zyklen wirkt auf jedem Element genauso wie \(\pi\). Damit ist die Zerlegung gezeigt.
\end{myproof}

\thm{Zerlegung in Transpositionen}{Jede Permutation \(\pi\in S_n\) mit \(n\geq2\) lässt sich als Produkt von Transpositionen schreiben.}

\begin{myproof}
Nach dem vorherigen Satz genügt es, einen Zyklus in Transpositionen zu zerlegen. Für \(m>1\) gilt
\[
(a_1\,a_2\,\cdots\,a_m)
=(a_1\,a_m)(a_1\,a_{m-1})\cdots(a_1\,a_2),
\]
wobei bei der Komposition von rechts nach links angewendet wird. Somit ist jeder Zyklus ein Produkt von Transpositionen und folglich jede Permutation.
\end{myproof}

\thm{Signum einer Permutation}{Für \(n\geq2\) gibt es genau eine Abbildung
\[
\operatorname{sgn}\colon S_n\to\{+1,-1\},
\]
für die gilt:
\begin{enumerate}[label=\alph*)]
\item Für jede Transposition \(\tau\) gilt \(\operatorname{sgn}(\tau)=-1\).
\item Für alle \(\pi,\sigma\in S_n\) gilt
\[
\operatorname{sgn}(\pi\sigma)
=\operatorname{sgn}(\pi)\operatorname{sgn}(\sigma).
\]
\end{enumerate}
Das Signum einer Permutation beschreibt, ob sie gerade oder ungerade ist.}

\begin{myproof}
\textbf{Eindeutigkeit.} Seien \(\operatorname{sgn}\) und \(\operatorname{sgn}'\) zwei Abbildungen mit den geforderten Eigenschaften. Nach dem vorherigen Satz lässt sich jede Permutation als Produkt von Transpositionen schreiben, etwa
\[
\pi=\tau_1\cdots\tau_k.
\]
Damit folgt
\[
\operatorname{sgn}(\pi)
=\operatorname{sgn}(\tau_1)\cdots\operatorname{sgn}(\tau_k)
=(-1)^k.
\]
Genauso gilt \(\operatorname{sgn}'(\pi)=(-1)^k\). Also stimmen die beiden Abbildungen überein.

\textbf{Existenz.} Wir definieren
\[
\operatorname{sgn}(\pi)
:=
\prod_{1\leq i<j\leq n}
\frac{\pi(j)-\pi(i)}{j-i}.
\]
Da \(\pi\) bijektiv ist, sind die Zahlen \(\pi(1),\ldots,\pi(n)\) paarweise verschieden. Die Faktoren sind daher alle von null verschieden. Die Produktformel ergibt stets \(+1\) oder \(-1\).

Durch Umordnen der Faktoren erhält man
\[
\operatorname{sgn}(\pi\sigma)
=\operatorname{sgn}(\pi)\operatorname{sgn}(\sigma).
\]
Für eine Transposition \(\tau=(1\,2)\) ergibt die Produktformel \(\operatorname{sgn}(\tau)=-1\). Jede Transposition \((i\,j)\) ist zu \((1\,2)\) durch Umbenennung der Elemente konjugiert; aus der Multiplikativität folgt daher ebenfalls \(\operatorname{sgn}((i\,j))=-1\).

Damit erfüllt die definierte Abbildung die beiden geforderten Eigenschaften. Zusammen mit der Eindeutigkeit ist der Satz bewiesen.
\end{myproof}

\cor{Parität der Anzahl von Transpositionen}{Jede Permutation kann entweder als Produkt einer geraden Anzahl oder als Produkt einer ungeraden Anzahl von Transpositionen geschrieben werden. Die beiden Möglichkeiten schließen sich gegenseitig aus.}

\begin{myproof}
Sei
\[
\pi=\tau_1\cdots\tau_k=\tau'_1\cdots\tau'_\ell
\]
eine Darstellung derselben Permutation als Produkt von Transpositionen. Wegen der Multiplikativität des Signums gilt
\[
(-1)^k=\operatorname{sgn}(\pi)=(-1)^\ell.
\]
Somit haben \(k\) und \(\ell\) dieselbe Parität: Sie sind beide gerade oder beide ungerade.
\end{myproof}

\dfn{Gerade und ungerade Permutationen}{Eine Permutation \(\pi\in S_n\) heißt \textbf{gerade}, wenn \(\operatorname{sgn}(\pi)=1\), und \textbf{ungerade}, wenn \(\operatorname{sgn}(\pi)=-1\).

Die Menge aller geraden Permutationen heißt \textbf{alternierende Gruppe}:
\[
A_n:=\{\pi\in S_n\mid \operatorname{sgn}(\pi)=1\}.
\]}

\thm{Die alternierende Gruppe ist eine Gruppe}{Die Menge \(A_n\) bildet bezüglich der Komposition von Permutationen eine Gruppe.}

\begin{myproof}
Die Identität ist gerade, also \(\operatorname{id}\in A_n\). Sind \(\pi,\sigma\in A_n\), so gilt
\[
\operatorname{sgn}(\pi\sigma)
=\operatorname{sgn}(\pi)\operatorname{sgn}(\sigma)
=1\cdot1=1.
\]
Also liegt \(\pi\sigma\in A_n\). Außerdem gilt
\[
1=\operatorname{sgn}(\operatorname{id})
=\operatorname{sgn}(\pi\pi^{-1})
=\operatorname{sgn}(\pi)\operatorname{sgn}(\pi^{-1}),
\]
und da \(\operatorname{sgn}(\pi)=1\) ist, folgt \(\operatorname{sgn}(\pi^{-1})=1\). Somit liegt auch \(\pi^{-1}\in A_n\).

Die Assoziativität der Komposition gilt bereits in \(S_n\) und damit auch in \(A_n\). Alle Gruppenaxiome sind erfüllt.
\end{myproof}

\mlemma{Ungerade Permutationen als Nebenmenge}{Sei \(\tau\in S_n\) eine Transposition. Dann ist
\[
A_n\tau:=\{\sigma\tau\mid\sigma\in A_n\}
\]
genau die Menge der ungeraden Permutationen.}

\begin{myproof}
Für \(\sigma\in A_n\) gilt
\[
\operatorname{sgn}(\sigma\tau)
=\operatorname{sgn}(\sigma)\operatorname{sgn}(\tau)
=1\cdot(-1)=-1.
\]
Also ist jedes Element von \(A_n\tau\) ungerade.

Sei umgekehrt \(\pi\) eine beliebige ungerade Permutation. Setze \(\sigma:=\pi\tau^{-1}\). Da \(\tau^{-1}=\tau\), gilt
\[
\operatorname{sgn}(\sigma)
=\operatorname{sgn}(\pi)\operatorname{sgn}(\tau^{-1})
=(-1)(-1)=1.
\]
Somit ist \(\sigma\in A_n\), und wegen \(\tau^2=\operatorname{id}\) gilt
\[
\pi=\sigma\tau.
\]
Also liegt jede ungerade Permutation in \(A_n\tau\).
\end{myproof}
